%%=============================================================================
%% De inhoud van je projectvoorstel
%%=============================================================================
%%
%% Zeven rubrieken, samen hoogstens twee bladzijden. Houd het kort: per
%% rubriek volstaan één tot drie zinnen. Een voorstel is geen rapport - het
%% moet je begeleider in vijf minuten laten beslissen of dit een geschikt
%% graduaatsproject is.
%%
%% Stem je voorstel af met je mentor VOORDAT je het indient. Je mentor
%% bevestigt dat dit is wat het bedrijf bedoelt en dat de tien werkdagen
%% haalbaar zijn.

%---------- 1. Context en probleemstelling ------------------------------------

\section{Context en probleemstelling}%
\label{vst:context}

% Waar loop je stage, wat doet dat team, en welk concreet probleem of welke
% concrete behoefte is de aanleiding? Noem de mensen voor wie het een probleem
% is. Vermijd algemeenheden als "bedrijven hebben nood aan".

Ik loop stage bij Belgateq, een bedrijf gespecialiseerd in keuringen. Bij het uitvoeren van keuringen wordt er een keuringsverslag ingevuld, maar als de keurder het keuringsverslag wilt aanpassen moet hij/zij steeds opnieuw van nul beginnen. 

%---------- 2. Doelstelling ---------------------------------------------------

\section{Doelstelling}%
\label{vst:doelstelling}

% Wat bestaat er op het einde van dit project dat er vandaag niet is? Formuleer
% het als een resultaat waarvan iemand kan vaststellen dat het er is, en
% vermeld waaraan je zal merken dat het geslaagd is.

Webpagina maken waarin keuringsverslag wordt ingevuld en later kan worden aangepast.

%---------- 3. Technologieën --------------------------------------------------

\section{Technologieën}%
\label{vst:technologieen}

% Met welke talen, frameworks, databanken en diensten ga je werken? Welke
% liggen vast door de omgeving van het bedrijf, en welke kies je zelf? Vermeld
% ook wat nieuw voor je is: technologische verbreding wordt gewaardeerd.
%
% Vermeld hier ook of je AI denkt in te zetten, en waarvoor. Dat mag en wordt
% zelfs verwacht; in je rapport verantwoord je later wat je precies gedaan hebt.
%
% Verwerkt je toepassing persoonsgegevens, of zit er AI in wat je bouwt? Zeg
% dat dan nu al in één zin: dan weet je begeleider meteen dat de AVG of de
% AI Act meespeelt, en kan dat mee bekeken worden bij de goedkeuring.

Er zal gebruikt gemaakt worden van JavaScript, css, HTML, PHP/Laravel, MySQL en AWS. PHP/Laravel en AWS is nieuw voor mij. AI zal gebruikt worden om sneller code te schrijven, te debuggen en tekst te vertalen.

%---------- 4. Afbakening -----------------------------------------------------

\section{Afbakening}%
\label{vst:afbakening}

% DE BELANGRIJKSTE RUBRIEK. Je stage en je graduaatsproject lopen bij hetzelfde
% bedrijf; dit is waar je ze uit elkaar houdt.
%
% Beschrijf:
%   - wat je dagelijkse stagewerk is;
%   - wat het afgebakende geheel is dat JIJ zelfstandig analyseert, ontwerpt
%     en oplevert - dat is je graduaatsproject;
%   - wat er uitdrukkelijk buiten valt.
%
% Toets voor jezelf: welke analyse en welk ontwerp maak jij zelf? Kan je daar
% niets over schrijven, dan heb je een taak en geen project.

Mijn dagelijkse stagewerk is het uitvoeren van kleine aanpassingen in de bestaande website. Bij mijn graduaatsproject zal ik een koppeling leggen tussen database en het formulier. Ik zal zelf het model en database ontwerpen na het analyseren van de bestaande code. De code blijft bij het bedrijf. Het valideren van de form in de frontend valt buiten de scope.

%---------- 5. Verwachte meerwaarde -------------------------------------------

\section{Verwachte meerwaarde}%
\label{vst:meerwaarde}

% Wat heeft het bedrijf hieraan? Wat blijft er van jouw werk over wanneer je
% stage afgelopen is? Denk aan tijdwinst, minder fouten, nieuwe mogelijkheden,
% lagere kosten, of een basis waarop verder gebouwd kan worden.

Keurders zullen niet meer van nul moeten beginnen wanneer ze een keuringsverslag willen aanpassen.

%---------- 6. Planning en werkbelasting --------------------------------------

\section{Planning en werkbelasting}%
\label{vst:planning}

% Het graduaatsproject is ongeveer tien werkdagen, samen ongeveer 100 uren.
% Verdeel die over de fasen van je project en geef per fase aan wat je oplevert.
% Vermeld ook of je die dagen gefocust opneemt of gespreid over de stageperiode,
% en dat je dit met je mentor afgesproken hebt.

\begin{table}[htbp]
  \centering
  \begin{tabular}{llr}
    \toprule
    \textbf{Fase} & \textbf{Wat lever je op} & \textbf{Uren} \\
    \midrule
    Analyse en vereisten      & overzicht van de vereisten   & 15  \\
    Technologiekeuze          & verantwoorde keuze           & 5  \\
    Ontwerp                   & architectuur en gegevensmodel& 20  \\
    Realisatie                & werkende oplossing           & 50  \\
    Rapport, poster en demo   & de op te leveren documenten  & 10  \\
    \midrule
    \textbf{Totaal}           &                              & \textbf{100} \\
    \bottomrule
  \end{tabular}
  \caption[Planning en werkbelasting]{\label{tab:vst-planning}%
    }
\end{table}

%---------- 7. Risico's -------------------------------------------------------

\section{Risico's}%
\label{vst:risicos}

% Wat kan dit project doen mislukken, en wat doe je dan? Denk aan: toegang tot
% gegevens of systemen die je nodig hebt, afhankelijkheid van een collega of
% een derde partij, technologie die je nog moet leren, of drukte op de
% werkvloer waardoor je dagen niet vrijkomen.
%
% Twee tot drie risico's met telkens één zin over wat je terugvalplan is.

Tijdsgebrek: Minder tijd steken in styling/pdf generatie.
Technologie die ik niet ken: s'avonds bijleren.
